\documentclass[11pt]{article}
\usepackage[margin=1in]{geometry}
\usepackage{amsmath,amssymb,amsthm,booktabs}
\usepackage[hidelinks]{hyperref}
\setlength{\parindent}{0pt}
\setlength{\parskip}{0.55em}
\newtheorem{proposition}{Proposition}
\newtheorem{corollary}{Corollary}
\newcommand{\argminset}{\operatorname*{arg\,min}}
\newcommand{\dd}{\mathrm{d}}
\setlength{\emergencystretch}{2em}
\title{S13 M2 Conditional Certificate Analysis\\\large Hidden lane allocation, endogenous receiving service, and common VSL}
\author{Methodology review artifact}
\date{30 September 2026}

\begin{document}
\maketitle

\section{Question and result}

S05 established two separate facts: lane allocation can be hidden by an exact
aggregate-history collision in a homogeneous conservative branch, and fixed
advice-independent service cannot convert such hidden information into a
different useful common speed action when fastest advice dominates arrivals.

S13 asks whether the admitted M2 class can bridge those facts without selecting
one favorable numerical coefficient. Under the model contract in the companion
document, the answer is conditionally positive:

\begin{quote}
There is a nonempty open parameter regime in which two states with identical
declared predecision aggregate histories have different unique optimal actions
over the same continuous lane-uniform VSL family. The difference arises because
advice changes preparation-zone residence time, which changes remaining
mandatory transfer and hence finite-receiving workload at the choke.
\end{quote}

This is a model-class existence certificate. It is not field validation,
calibration, or proof that real lane contrast is necessary for VSL control.

\section{Contract recalled}

Let total cohort mass be \(M>0\). Hidden state \(s\in\{L,H\}\) has ending-lane
fraction \(p_s\), with \(0<p_L<p_H<1\). Both states have the same aggregate
predecision history and total mass. One lane-uniform advice speed \(v\) acts
through preparation-zone traversal time
\[
 \tau=L/v\in[\tau_{\min},\tau_{\max}].
\]
On a fixed continuous mandatory-transfer branch, ending-lane mass remaining at
the choke is
\[
 e_s(\tau)=p_sM e^{-\lambda\tau},\qquad \lambda>0.
\]
Remaining mandatory-transfer mass consumes receiving-work factor \(\alpha>1\)
at a choke with receiving-work rate \(\mu>0\). The complete per-mass time score is
\begin{equation}
 J_s(\tau)=\tau+\frac{M}{2\mu}+Kp_s e^{-\lambda\tau},
 \qquad K=\frac{(\alpha-1)M}{2\mu}>0.
 \label{eq:J}
\end{equation}

The first term is full preparation delay. The constant is ordinary cohort
clearance. The last term is additional finite-receiving time caused by mandatory
merging still present at the choke. Thus slower advice pays travel time but can
reduce service workload.

\section{Exact observer handoff}

\begin{proposition}[Conditional exact predecision collision]
Suppose the S05 homogeneous conservative preparation reaches the entrance of the
fixed M2 preparation zone at the decision time. If lane transfers before that
time preserve aggregate mass, cumulative counts, and common mean speed, then two
states with ending-lane fractions \(p_L\ne p_H\) have identical declared strong
aggregate histories at the decision boundary.
\end{proposition}

\begin{proof}
The S05 construction gives pointwise equality of every declared aggregate
history channel while allowing lane-transfer terms to cancel in the lane sum.
Choose the spatial entrance to the M2 zone as the decision boundary and require
the common cohort front to reach it first at that time. The roadway geometry and
M2 zone exist in both cases; only the already-hidden lane decomposition differs.
No M2 postdecision evolution has yet occurred, so it cannot leak information
into the predecision history. Therefore the observer histories are equal while
\(p_L\ne p_H\). This proposition remains conditional on the S05 reachability
assumptions; it does not certify a native simulator trajectory.
\end{proof}

\section{Action-split certificate}

Differentiate equation~\eqref{eq:J} with respect to traversal time:
\begin{align}
 J_s'(\tau)&=1-K\lambda p_s e^{-\lambda\tau},
 \label{eq:first}\\
 J_s''(\tau)&=K\lambda^2p_s e^{-\lambda\tau}>0.
 \label{eq:second}
\end{align}
Each state therefore has a strictly convex score over the full action interval.

\begin{proposition}[Different unique useful common actions]
If
\begin{equation}
 K\lambda p_L<e^{\lambda\tau_{\min}}
 <K\lambda p_H<e^{\lambda\tau_{\max}},
 \label{eq:regime}
\end{equation}
then the low-ending-lane state uniquely selects fastest advice,
\begin{equation}
 \argminset_{\tau\in[\tau_{\min},\tau_{\max}]}J_L(\tau)
 =\{\tau_{\min}\},
 \label{eq:low}
\end{equation}
while the high-ending-lane state uniquely selects the slower interior action
\begin{equation}
 \argminset_{\tau\in[\tau_{\min},\tau_{\max}]}J_H(\tau)
 =\left\{\tau_H^*\right\},\qquad
 \tau_H^*=\frac{1}{\lambda}\log(K\lambda p_H)>
 \tau_{\min}.
 \label{eq:high}
\end{equation}
The optimal-action sets are disjoint.
\end{proposition}

\begin{proof}
For state \(L\), the first inequality in~\eqref{eq:regime} gives
\(J_L'(\tau_{\min})>0\). Equation~\eqref{eq:second} shows that the derivative is
strictly increasing, so \(J_L'(\tau)>0\) throughout the interval. The unique
minimum is \(\tau_{\min}\), equivalent to \(v_{\max}\).

For state \(H\), the middle inequality gives
\(J_H'(\tau_{\min})<0\), while the final inequality gives
\(J_H'(\tau_{\max})>0\). Strict convexity yields one zero of the derivative in
the interior. Solving equation~\eqref{eq:first} gives~\eqref{eq:high}. Since
\(\tau_H^*>\tau_{\min}\), the action sets are disjoint.
\end{proof}

The inequalities are strict, so the result persists under small parameter
perturbations. It is therefore an open-regime certificate rather than a
one-point favorable example. A normalized illustration such as
\(\lambda\tau_{\min}=0.2\), \(\lambda\tau_{\max}=2\),
\(K\lambda=4\), \(p_L=0.2\), and \(p_H=0.5\) satisfies the inequalities; these
numbers demonstrate nonemptiness and are not traffic calibration.

\section{Why the S05 obstruction no longer applies}

In S05, advice changed cumulative arrivals while each hidden state's service
process remained fixed with respect to advice. Here service workload is
\[
 W_s(\tau)=M+(\alpha-1)p_sM e^{-\lambda\tau},
\]
so the action changes service through remaining mandatory transfer. Faster
advice shortens travel but leaves more ending-lane mass to consume scarce
receiving work. Slower advice can reduce that work. The arrival-only monotonicity
premise is therefore false inside the M2 regime.

The certificate does not contradict work conservation. Lateral transfer is
paired, all mass clears, and the action-induced travel time is fully charged.
The benefit arises from changing receiving work, not moving uncounted waiting
upstream.

\section{Observer-only aggregate decision limit}

\begin{proposition}[Positive observer-only regret]
At the decision boundary, any deterministic controller whose input is limited to
the declared strong aggregate history must select the same action for states
\(L\) and \(H\). Under~\eqref{eq:regime}, no one action is optimal for both, and
the minimum worst-state regret is strictly positive.
\end{proposition}

\begin{proof}
The observer histories and feasible action family are identical, so a
deterministic observer-only rule returns the same \(\tau\) in both states. Define
\[
 R_s(\tau)=J_s(\tau)-\min_x J_s(x)\ge0.
\]
Strict convexity implies \(R_s(\tau)=0\) only at state \(s\)'s unique minimizer.
The two zero sets are disjoint. Hence the continuous function
\(\max\{R_L(\tau),R_H(\tau)\}\) is positive everywhere on the compact action
interval and has a strictly positive minimum.
\end{proof}

A source-backed deterministic aggregate capacity-drop model receiving the same
history and action likewise produces one prediction at the decision point. It
can match one hidden case, an average, or later observed divergence, but it
cannot condition the initial action on \(p_s\) without adding information outside
the declared observer. This is an information result, not a claim that every
possible feedback controller remains blind after postdecision measurements
diverge.

\section{Conditional gradient path}

The action-to-outcome chain is
\[
 v\longrightarrow \tau=L/v
 \longrightarrow e_s=p_sM e^{-\lambda\tau}
 \longrightarrow W_s=M+(\alpha-1)e_s
 \longrightarrow J_s=\tau+W_s/(2\mu).
\]
Inside the fixed active regime,
\begin{equation}
 \frac{\dd J_s}{\dd v}
 =-\frac{L}{v^2}\left(1-K\lambda p_s e^{-\lambda L/v}\right).
 \label{eq:gradient}
\end{equation}

\begin{corollary}[State-relevant gradient direction]
Under~\eqref{eq:regime}, at fastest advice the low-ending-lane state has
\(\dd J_L/\dd v<0\): increasing speed toward the endpoint reduces cost. The
high-ending-lane state has \(\dd J_H/\dd v>0\): reducing speed from the endpoint
reduces cost. Thus the hidden state changes the locally useful gradient direction
inside the conditional model.
\end{corollary}

This establishes conceptual downstream relevance in the analytical member, not
automatic-differentiation correctness for a future implementation. At receiving,
demand/supply, zero-mass, endpoint, or priority active-set changes, the smooth
derivative may need one-sided or piecewise treatment.

\section{Status against the core review questions}

\begin{center}
\begin{tabular}{p{0.33\textwidth}p{0.57\textwidth}}
\toprule
Question & Conditional S13 answer \\
\midrule
What hidden macro-omitted information can alter control? & Ending-lane fraction,
because it determines remaining mandatory-transfer receiving work while the
declared aggregate history is identical. \\
Can it alter the controller result? & Yes within the open M2 regime: the two
states have different unique optimal actions and positive observer-only minimax
regret. \\
Can it provide relevant gradient information? & Yes conceptually within a fixed
active regime: the hidden state can reverse the locally useful gradient direction
at fastest advice. \\
Is lane contrast empirically necessary? & Not established. The contract is a
conditional analytical model, not calibrated or cross-fidelity evidence. \\
\bottomrule
\end{tabular}
\end{center}

\section{Falsification and limits}

The positive result disappears in aggregate/fixed-service limits
\(\alpha\to1\), \(\lambda\to0\), equal hidden fractions, or nonbinding receiving.
It also fails outside~\eqref{eq:regime}, where both states may choose fastest
advice or the same interior action.

The decisive unsupported empirical join remains the constant positive transfer
branch under common VSL. It is source-compatible but not calibrated. The batch
and homogeneous proportional service assumptions remove order and transient
wave detail. The exact observer collision is inherited conditionally from S05,
not demonstrated in native traffic. Therefore the result establishes
methodological feasibility and a complete conditional chain, not field truth or
minimum mesoscopic state.

\section{Verdict}

Under the approved S13 contracts, the result is
\textbf{M2C-A---conditional decision-relevant certificate}. All six planned
analytical conditions pass inside a nonempty open parameter regime. The
remaining gap is empirical/model-validity support for the selected response
closure, followed by implementation-level derivative correctness only if a
future project chooses to build it.

\end{document}
